% SOURCE_AUTHORITY: sga5_fr_workpass.tex, lines 12497--12509 (LF-normalized unit).
% SOURCE_SHA256: 791F4EFFC5E02832D5D77ED03518C8156D6F07E4C8238B03545DB93D883FBB28
Com efeito, graças à hipótese $P^\bullet\in D(\Lambda[G])_{\parf}$, pode-se supor que $P^\bullet$ é limitado e que todas as suas componentes são projetivas de tipo finito. Tem-se, portanto,
\[
\RHom_{\Lambda[G]}(P^\bullet,M^\bullet)=\Hom_{\Lambda[G]}^\bullet(P^\bullet,M^\bullet)
\]
% SOURCE_ANOMALY_PRESERVED: the frozen authority and scan print reference (4.3.4).
e podem-se utilizar os isomorfismos (4.3.4). Mas $\check P^\bullet$, assim como $P^\bullet$, é limitado e de componentes projetivas de tipo finito, de modo que
\[
\check P^\bullet\otimes_{\Lambda[G]}M^\bullet\simeq \check P^\bullet\otimes_{\Lambda[G]}^{\mathbf L}M^\bullet.
\]
Da mesma forma, $\check P^\bullet\otimes_\Lambda M^\bullet=\check P^\bullet\otimes_\Lambda^{\mathbf L}M^\bullet$. Além disso, os componentes do complexo $\check P^\bullet\otimes_\Lambda M^\bullet$ são cohomologicamente triviais ([3] IX 3), e tem, portanto:
\[
\R\Gamma^G(\check P^\bullet\otimes_\Lambda^{\mathbf L}M^\bullet)\simeq (\check P^\bullet\otimes_\Lambda M^\bullet)^G.
\]
